\subsection{All weights in the class-two quotient}

The next result is stronger in weight range than the finite
calculation but concerns the quotient by rational generator
length at least three. Its induced filtration need not be the
quotient of the full arithmetic filtration.

\begin{theorem}[K17-2-CLASS2]\label{thm:2-class2}
Let $H^{(2)}$ be the actual compact image of $H$ in the rational
class-two quotient of the free odd Lie algebra, with its induced
arithmetic-weight filtration. Its integral graded Lie algebra
is the free class-two $\mathbb Z_2$ Lie algebra on one primitive
generator in each odd weight $m_i=2i+3$. The primitive coefficient
in the canonical coordinate $s_{m_i}$ is, up to a unit,
\[
 \delta_i=\prod_{j<i}(5^{m_i}-5^{m_j}),\qquad
 v_2(\delta_i)=3i+v_2(i!).
\]
\end{theorem}

\begin{proof}
Fix $r\geq1$ and cutoff $W=2r+2$. Put $a_i=5^{2i+3}$ for
$0\leq i<r$, and
\[
 P_i(X)=\prod_{j<i}(X-a_j),\qquad
 P_r(X)=\sum_{t=0}^r c_tX^t.
\]
The linear image $A$ of $H$ is the evaluation image of
$\Lambda$ at these $r$ nodes, using $\gamma$ as variable.
Successive monic division gives kernel $P_r\Lambda$.
The classes of $1,\gamma,\ldots,\gamma^{r-1}$ and of
$P_0,\ldots,P_{r-1}$ are two integral bases. Both $\gamma$
and its inverse preserve $A$, since $P_r$ has unit constant
coefficient.

The image $C$ of the closed commutator subgroup is central
and is the image of $\bigwedge^2 A$. Its basis vectors are
\[
 P_i(a_k)P_j(a_l)-P_j(a_k)P_i(a_l),\qquad
       i<j,\quad i+j\leq r-2,
\]
in the target coordinates $[s_{2k+3},s_{2l+3}]$.
They are triangular by increasing $i+j$, with diagonal
$\delta_i\delta_j$: a nonzero evaluation requires
$k+l\geq i+j$, and equality forces the same pair.
Pairs with $i+j>r-2$ vanish at this cutoff. These are actual
commutators of lifts of the Newton basis, and $C$ is stable
under $\gamma^{\pm1}$ by its exterior-square description.

Set $h_t=\gamma^t h\gamma^{-t}$ and form the ordered word
\[
 R_r=\prod_{t=0}^r h_t^{c_t}.
\]
Its linear image is zero. Let $E_r$ be its central logarithm
at the cutoff. The kernel of the linear projection of the
compact image is generated by $C$ and the $\gamma$-orbit of
$E_r$: modulo the closed commutator subgroup the relevant
classes are $P_r\Lambda$, generated by the conjugates of
$R_r$, and compactness supplies their closed span. Thus
$E_r\in C$ would identify the whole kernel with $C$.
Here ``central'' refers to the length-two derived subspace;
other central length-one directions at the cutoff are not
part of this kernel.

For nonnegative integral $t$, the coefficient of
$[s_m,s_n]$ in $\log h_t$ is the polynomial
\[
 b_t(X,Y)=-\frac12\bigl(Y^t S_t(X)-X^tS_t(Y)\bigr),
 \qquad S_t(X)=\sum_{j=0}^{t-1}X^j,\quad S_0=0,
\]
evaluated at $X=5^m$, $Y=5^n$. This follows by summing the
geometric series for
$D=\operatorname{Ad}(\exp(-a/2))\tau$ in class two.
The exact class-two product formula therefore gives
\[
 Z_r(X,Y)=\sum_t c_tb_t(X,Y)
 +\frac12\sum_{x<y}c_xc_y(X^xY^y-X^yY^x)               \tag{C2.1}
\]
for the central coefficient of $R_r$. This alternating
polynomial belongs to $\tfrac12\mathbb Z[X,Y]$ and has degree
at most $r$ in each variable. Monic division in the tensor
Newton basis gives
\[
 Z_r=\sum_{i<j}z_{ij}
       \bigl(P_i(X)P_j(Y)-P_j(X)P_i(Y)\bigr),
       \qquad z_{ij}\in\tfrac12\mathbb Z.             \tag{C2.2}
\]

We prove $z_{ij}\in\mathbb Z$ for $i+j\leq r-2$ by parity.
All nodes are odd. In $\mathbb F_2[X,Y]$, put
$A=X-1$, $B=Y-1$, and $Q=A+B+AB=XY-1$.
Then $P_r(X)=A^r$, $P_r(Y)=B^r$, $P_r(XY)=Q^r$.
The linear part of twice (C2.1) is
$(Q^r+A^r)/B+(Q^r+B^r)/A$, with polynomial quotients.
The cross term is $P_r(X)P_r(Y)+P_r(XY)$ because
$c_t^2=c_t$ modulo $2$. Hence
\[
 2Z_r=\frac{Q^r+A^r}{B}+\frac{Q^r+B^r}{A}
                      +A^rB^r+Q^r
       \in(A,B)^{r-1}\quad\text{over }\mathbb F_2.    \tag{C2.3}
\]
Every quotient displayed here has minimum total degree
$r-1$. Since $P_i(X)P_j(Y)$ reduces to $A^iB^j$, the
coefficient $2z_{ij}$ is even for $i+j<r-1$. Its known
half-integrality then proves the required integrality.

All remaining tensor Newton terms vanish at the cutoff.
Equation (C2.2) therefore puts $E_r$ in the actual lattice
$C$, and its whole $\gamma$-orbit stays there. This proves
that the full linear kernel is exactly $C$.

Divide $R_r$ by an actual product of commutators with the
same central image through $2r+2$. The new word has unchanged
abelian class $P_r$ and no class-two component below weight
$2r+3$; its next linear coefficient is $\delta_r$. Thus every
odd Newton minimum is attained in $H^{(2)}$. At even weights,
the just identified lattice $C$ has the leading basis
$\delta_i\delta_j[s_{m_i},s_{m_j}]$ by triangularity. At odd
weights there is no length-two component. These observations
prove the theorem.
\end{proof}

The exact script \path{evidence/work/pro2_integral_class2_newton.py}
checks the group-product formula against its geometric-series
form and eliminates in the Newton-wedge basis. Its finite
checks helped identify (C2.3); the polynomial identity proves
the all-weight assertion. The result does not assert that
these corrected words have no components of generator length
at least three below their intended odd weight in the full
image. The higher-length lifting is supplied by the separate induction
in Theorem~\ref{thm:2-allweights}.
